% 第二章 正文文字格式
\thesischapterpage{2\quad 正文文字格式}

\section{正文文字格式}

\subsection{论文正文}
\thesisbodyfont
论文正文是主体，一般由标题、文字叙述、图、表格和公式等部分构成。一般可包括理论分析、计算方法、实验装置和测试方法，经过整理加工的实验结果分析和讨论，与理论计算结果的比较以及本研究方法与已有研究方法的比较等，因学科性质不同可有所变化。\par
正文内容应层次清楚、逻辑严密、语言准确。章节标题应按论文结构逐级组织，正文段落一般采用首行缩进，图表和公式应在正文中进行必要说明。

\subsection{字数要求}
\subsubsection{本科论文字数要求}
\thesisbodyfont
论文主体部分字数要求：理工类专业一般不少于1.5万字，其他专业一般不少于1.0万字。

\subsection{本章小结}
\thesisbodyfont
本章给出了正文文字排版的示例。正式写作时，可保留本文件结构并替换为对应章节内容。
